\chapter{Background Gauge Fields and Index Theory}
\label{ch:curved-background-gauge-fields-index-theory}

\ConventionsEuclidean

Background gauge fields couple QFT to vector bundles over spacetime.  For
fermions, the index of the coupled Dirac operator controls chiral zero modes,
global charge violation in fixed backgrounds, and the topological part of
anomaly inflow.

\section{Dirac Operators with Background Bundles}

Let \(M^{2n}\) be a closed oriented Riemannian spin manifold with spinor
bundles \(S^\pm\).  Let \(E\to M\) be a Hermitian vector bundle with unitary
connection \(A\).  The chiral Dirac operator is
\[
  D_A^+:\Gamma(S^+\otimes E)\longrightarrow \Gamma(S^-\otimes E).
\]
Its Fredholm index is
\[
  \operatorname{Ind}(D_A^+)
  =
  \dim\ker D_A^+-\dim\ker D_A^- .
\]
The heat-kernel representative of the index is
\[
  \operatorname{Ind}(D_A^+)
  =
  \lim_{t\to0^+}\operatorname{Tr}\!\left(
  \Gamma\,e^{-tD_A^2}
  \right),
\]
where \(\Gamma\) is chirality on \(S^+\oplus S^-\).

\section{Atiyah--Singer Formula}

The index theorem gives
\[
  \operatorname{Ind}(D_A^+)
  =
  \int_M \widehat A(TM)\operatorname{ch}(E).
\]
Here
\[
  \operatorname{ch}(E)=\operatorname{tr}_E\exp\!\left(\frac{\ii F_A}{2\pi}\right),
\]
and
\[
  \widehat A(TM)=1-\frac{p_1(TM)}{24}+\cdots
\]
is the \(\widehat A\)-class of the tangent bundle.  In four dimensions,
\[
  \operatorname{Ind}(D_A^+)
  =
  \int_M
  \left[
  \operatorname{ch}_2(E)
  -\frac{\operatorname{rank}(E)}{24}p_1(TM)
  \right].
\]
The trace in \(\operatorname{ch}(E)\) is the trace in the fermion
representation.  Converting this formula to the gauge-theory trace convention
requires the representation index relating the two bilinear forms.

\section{Zero Modes and Charge Violation}

Let a left-handed Weyl fermion of charge representation \(E\) be placed in a
background \(A\).  The Berezin integral over fermions vanishes unless
insertions absorb the zero modes of \(D_A^+\).  If
\[
  k=\operatorname{Ind}(D_A^+)>0
\]
and there are no right-handed zero modes in the conjugate sector, correlation
functions must contain \(k\) insertions of the corresponding fermionic
variable to be nonzero.  This is an exact finite-dimensional statement about
the zero-mode part of the Berezin integral after the nonzero modes have been
paired.

\section{Anomaly Polynomial}

For a chiral fermion in \(2n\) dimensions, the anomaly polynomial is the
\((2n+2)\)-form component
\[
  I_{2n+2}
  =
  \left[\widehat A(TM)\operatorname{ch}(E)\right]_{2n+2}.
\]
Choose a Chern-Simons form \(I_{2n+1}\) with
\[
  dI_{2n+1}=I_{2n+2}.
\]
The descent equation
\[
  \delta_\alpha I_{2n+1}=dI_{2n}^{(1)}(\alpha)
\]
gives the local anomalous variation
\[
  \delta_\alpha W[A,g]
  =
  2\pi\ii\int_M I_{2n}^{(1)}(\alpha).
\]
This formula agrees with the inflow construction of
Chapter~\ref{ch:global-anomaly-inflow-invertible-field-theories} when the
bulk invertible theory is defined by the corresponding differential
cohomology refinement of \(I_{2n+2}\).

\section{Four-Dimensional Gauge and Mixed Terms}

In four dimensions the six-form component contains
\[
  I_6
  =
  \left[\operatorname{ch}_3(E)
  -\frac{p_1(TM)}{24}\operatorname{ch}_1(E)\right].
\]
The first term gives the cubic gauge anomaly.  The second term gives the
mixed gauge-gravitational anomaly for a \(U(1)\) factor.  For a semisimple
gauge group, \(\operatorname{ch}_1(E)=0\), so this mixed term is absent in
that representation.  The condition for gauging a background symmetry is the
vanishing of the total anomaly polynomial modulo terms that are canceled by
specified inflow or by an explicitly declared local counterterm.

\section{Families and Global Anomalies}

Let \(\mathcal A/\mathcal G\) be a background-field moduli space.  The
fermion determinant is a section of a determinant line over this space:
\[
  \operatorname{Det} D
  =
  \bigwedge^{\max}\ker D
  \otimes
  \left(\bigwedge^{\max}\operatorname{coker}D\right)^\ast .
\]
Local anomalies are the curvature of this line with its natural connection.
Global anomalies are its holonomies around loops in \(\mathcal A/\mathcal G\).
Trivializing the determinant line, not merely its curvature, is required for a
single-valued background-field partition function.

\begin{openproblem}[Index-theoretic anomaly line in interacting QFT]
\label{op:index-anomaly-line-interacting-qft}
Extend the determinant-line and index-theoretic description from free chiral
fermions coupled to backgrounds to interacting continuum QFTs with composite
currents, contact terms, and extended operators, while retaining an explicit
comparison with local anomaly descent.
\end{openproblem}
